\documentclass[letterpaper]{article}
\usepackage{natbib}
\usepackage[utf8]{inputenc}
\usepackage[margin=3.5cm]{geometry}
\usepackage{listings}
\usepackage{adjustbox}
\usepackage{xcolor}
\usepackage{verbatim}
\usepackage{graphicx}%images
\usepackage{fancyhdr}%for headers and footers
\usepackage{adjustbox}
\usepackage{verbatim}
\usepackage{listings}
\usepackage{fancyhdr}
\usepackage{multirow}
\usepackage{amsmath}
\usepackage{array}
\usepackage{mathpazo}
\usepackage{subcaption}
\usepackage{float}
\usepackage{csvsimple}
\usepackage{filecontents}
\usepackage{lscape}
\usepackage{afterpage}
\usepackage{hyperref}
\usepackage{inconsolata}
\usepackage{color}
\usepackage{todonotes}

\newenvironment{conditions}
  {\par\vspace{\abovedisplayskip}\noindent\begin{tabular}
  {>{$}l<{$} @{${}={}$} l}}
  {\end{tabular}\par\vspace{\belowdisplayskip}}


  \newenvironment{conditions_bis}
  {\par\vspace{\abovedisplayskip}\noindent\begin{tabular}
  {>{$}l<{$} @{${}\in{}$} l}}
  {\end{tabular}\par\vspace{\belowdisplayskip}}

\definecolor{pblue}{rgb}{0.13,0.13,1}
\definecolor{pgreen}{rgb}{0,0.5,0}
\definecolor{pred}{rgb}{0.9,0,0}
\definecolor{pgrey}{rgb}{0.46,0.45,0.48}

\lstset{language=Python,
  showspaces=false,
  numbers=left,
  showtabs=false,
  breaklines=true,
  showstringspaces=false,
  breakatwhitespace=true,
  commentstyle=\color{pgreen},
  keywordstyle=\color{pblue},
  stringstyle=\color{pred},
  basicstyle=\ttfamily,
  moredelim=[il][\textcolor{pgrey}]{$ $},
  moredelim=[is][\textcolor{pgrey}]{\%\%}{\%\%}
}



\hypersetup{
    colorlinks,
    citecolor=black,
    filecolor=black,
    linkcolor=black,
    urlcolor=black
}


% ------ HEADERS AND FOOTERS -----------
% \lhead{INFO-F403}
% \rhead{Project Report - Part 1}
%\pagestyle{fancy}
% \rfoot{\thepage}
%\cfoot{}
%\lfoot{Academic Year 2017-18}

\newcommand{\HRule}{\rule{\linewidth}{0.5mm}} %newcommand for cover page

\begin{document}

\begin{titlepage}
\begin{center}


\textsc{\LARGE universit\'e libre de bruxelles}\\[1.0cm]
\textsc{\Large D\'epartment d'Informatique}\\[1.5cm]

% Upper part of the page. The '~' is needed because \\
% only works if a paragraph has started.
\includegraphics[width=0.3\textwidth]{images/ulblogo.jpg}~\\[1cm]

\textsc{
\large MEMO-F-403 \\
\Large  Preparatory work for the master thesis
 \\[1cm]}
% Title
\HRule \\[0.7cm]

{ \huge \bfseries Thesis - draft \#1 \\[0.7cm] }

\HRule \\[2cm]

% Author and supervisor
\noindent
\begin{center} \large

\emph{Author:}\\
\Large Hakim \textsc{Boulahya}\\
\end{center}
\begin{center} \large

\emph{Supervisor:} \\
Emmanuel \textsc{Filiot} \\

\end{center}

\vfill

% Bottom of the page
{\large \today}

\end{center}
\end{titlepage}

\tableofcontents
\newpage

\section{Introduction}


\subsection{Motivation}

\paragraph{}


Automata theory is used in various field in Computer Science, especially
in Computer-Aided Verification. Multiple algorithms has been theorized
to provide more efficient way to resolve those problem
\todo{Talk about
the problems, complexity and alternative (Safra vs antichain)}.


\paragraph{}

\todo{Talk about unversality and LTL real.}
Those different algorithms used a data structure called antichain.
An antichain is a subset of a set that allow to represent it, in a more
compact way.

\subsection{Objective}

The goal of this thesis is to provide an efficient implementation of
data structures used in those algorithms, especially antichain. We will
mainly focus on implementing antichain-related data structure and provide
a library to be used in different tool such as Acacia+, Owl.

The first, and main, objective is to implement an Antichain object
in Java to be used in Owl. Then if possible used this implementation
for other tool such as Acacia+, by either providing bindings or other.

As we mainly focus on effiency, it could be interesting to use a C
implementation and provide binding to a Java class.
\todo{This small paragraph is an open discussion}

\subsection{Related work}

\paragraph{}

AaPAL library is a that was implemented in the context of Bohy's PhD thesis
to provide an antichain library and be able to implement the antichains
based algorithm for the synthesis problem.

\paragraph{}

Java already provide built-in implementation for Set.

\section{Preliminaries}

In this section, we will provide formal definitions of the data
structures that we will implement. We recall the notion of set, partially
ordered set and total order set. We then give a definition for antichains.

The definitions and examples for this section are based on ...\todo{insert
ref of examples}

\subsection{Formal definitions}


\paragraph{Binary relation}

A binary relation for an arbitrary set $S$ is
a set of pair $R \subseteq S \times S$.
There are four important properties: reflexitivity, transitivity,
symmetry, antisymmetry and total.

\paragraph{}

A relation R on S is said to be:

\begin{itemize}
    \item Reflexive:
    iff for all $s \in S$ it holds that $(s, s) \in R$
    \item Transitive:
    iff for all $s_1, s_2, s_3 \in S$,
    if ($s1, s2) \in R$ and $(s2, s3) \in R$
    then it holds that $(s1, s3) \in R$
    \item Symmetric: iff $(s_1, s_2) \in R$ then $(s_2, s_1) \in R$.
    \item Antisymmetric: iff $(s_1, s_2) \in R$
    and $(s_2, s_1) \in R$ then $s_1 = s_2$
    \item Total: \todo{total definition}
\end{itemize}

\paragraph{Order}

A \textit{partial order} is a binary relation that is reflexive,
transitive and antisymmetric. We note a
partial order relation by ... To note the appartenance of a binary
relation to a partial order, we note s1 .. s2 which is equivalent
to (s1, s2) in ...
A total order is a partial order that is total.
\todo{give definition of (un)comparable using partial order definition}

\todo{Is this total order affirmation correct ?}
In this thesis we are more interested in partial ordered sets as
total order sets can be easily implemented as lists.

\paragraph{Ordered sets}

Partial/Order sets

Lattice (semi upper and lower)

Closed sets

Antichain

Includes propreties to implement

\subsection{Existing implementation}


\todo{Everything below is a fast/draft notes}

\section{Motivation and objective}

For the moment, the implementation to represent the data strutures in
 Acacia+
is specifically designed for a specific set. The idea is to propose a new
library implementation to provide an API that will implement important
data structures that are used in syntesis algorithms.

\paragraph{}

The objective of the thesis is to provide an efficient library to represent
antichain data structures and implement different operation. An final
goal is to be able to use this library within Aaron/Acacia.


* Impl. datastrcture antichain in Java
* Theoritacal context: synthesis,
unverversality and automata theory known problem
* Practical context: Owl, Acacia+ and AaPAL

* Why interesting: More efficient than CTL symbolic with BDD

\subsection{Related Work}

* Java Built-in IMPL
* AaPAL
* Stackoverflow
* Acacia ?

\section{Notions}

\subsection{Model checking and synthesis}

Model checking is the process to verify a software model with
specifications.

Synthesis is the process to derive the system from specifications.

\cite{bohyphd}

\cite{ltl_rea}
\subsection{Data structures}

\paragraph{Binary Decision Diagram}

\paragraph{Antichain}

BDD vs antichain

Antichain vs pseudo-antichain

\section{Questions}

* Maastricht library ? \cite{acacia} \cite{aapal}
* How to references ?
* What to impl (about operations) ?

* Is implementing LTL Rea. or Universality a final goal of the thesis ?

\bibliographystyle{alpha}
\bibliography{thesis}

\end{document}
